\section{Related Work}

\noindent \textbf{Bi-modal Embedding Models.}
Early vision--language embedding models primarily focus on aligning text and images in a shared space, typically via dual-encoder contrastive learning~\citep{ALIGN,CLIP,BLIP,blip2}.
For example, CLIP~\citep{CLIP} learns image and text encoders by matching paired image--caption data at scale.
More recently, a line of work builds retrieval-oriented bi-encoders on top of modern VLMs~\citep{MMEB,GME,mme5,B3,Moca}, converting generative or instruction-following backbones into embedding models.
For example, mmE5~\citep{mme5} improves text--image embedding quality through data synthesis.
Despite strong progress on text--image representations, these bi-modal models are not designed to fully meet modern retrieval needs that increasingly involve additional modalities.

\noindent \textbf{Omni-modal Embedding Models.}
Moving beyond text--image, recent work has started to incorporate video~\citep{MMEB-v2,UME-R1,UniME-V2}.
For example, UME-R1~\citep{UME-R1} explores a reasoning-driven generative embedding paradigm that leverages reinforcement learning.
More recent efforts further expand to video and audio modalities~\citep{tevatron-omni,Nemotron,LCO-Embed}.
For example, Omni-Embed-Nemotron~\citep{Nemotron} builds a unified embedding model on a Qwen-Omni backbone to support omni-modal retrieval in a single shared space.
Similarly, LCO-Embed~\citep{LCO-Embed} proposes a language-centric omni-modal embedding framework.

While these omni-modal models broaden coverage, they largely rely on VLM pretraining to provide implicit alignment and typically lack explicit calibration and alignment mechanisms.
This limitation becomes more critical as the number of supported modalities grows.
Unlike these models, \ours{} introduces a lightweight, plug-and-play explicit alignment recipe to improve robustness under mixed-modality training.



